\section{Discussion}
\label{sec:discussion}

\paragraph{What the guarantee adds.}
A direct certificate gives each exclusion a specific source and start time.
The local-influence theorem shows how far a changed replacement decision can propagate.
The prefix theorem additionally controls which periods it can affect.
These properties support debugging and incremental maintenance.
An operator can inspect the witness responsible for a missing passage without reconstructing an entire component.
With a fixed candidate ranking, a single corrected pair requires reconsidering one eligibility bit at the query cutoff.

\paragraph{Why accurate grouping can still be sufficient.}
Exact-group date selection is a strong baseline when extracted keys remain stable.
Its historical TempLAMA accuracy matches or exceeds the direct construction in several comparisons.
Direct certificates become useful when pair decisions carry information that exact equality misses.
They preserve that information without assuming transitivity.
The RealProse result illustrates the benefit of approximate matching.
The prefix and sensitivity results identify the additional cost of converting those matches into components.

\paragraph{From replacement to coexistence.}
The TempLAMA source audit reveals that answer order cannot establish replacement.
A previous employer, team, or office can remain valid when another value appears.
Complete answer sets expose this distinction during evaluation.
A deployment gate should also express it during inference.
For instance, evidence of a new concurrent role should not automatically close an existing role.
The compiler can accept such a gate unchanged.
Its formal guarantees then apply to that gate's decisions.

\paragraph{Practical integration.}
The certificate layer sits between relevance ranking and context selection.
It can therefore reuse an existing retriever and reader.
The source text remains available for citation and historical retrieval.
When an application's metadata supplies actual effective dates, those dates define the applicable periods.
When it supplies publication dates, the certificates instead describe the policy over those dated records.
The application must choose that interpretation before compiling its memory.
The same choice determines what a dated user question asks the system to retrieve.
